\documentclass[a4paper,12pt]{article}
\usepackage[utf8]{inputenc}
\usepackage[T1]{fontenc}
\usepackage[ngerman]{babel}
\usepackage{geometry}
\geometry{margin=2.5cm}
\usepackage{amsmath}
\usepackage{booktabs}
\usepackage{hyperref}
\hypersetup{colorlinks=true, linkcolor=blue, urlcolor=blue}

\begin{document}

% =========================================================
\section*{Discretization-invariance? On the Discretization Mismatch
          Errors in Neural Operators}
% =========================================================

\noindent\textbf{Autoren:} Wenhan Gao, Ruichen Xu, Yuefan Deng, Yi Liu\\[2pt]
\textbf{Institution:} Department of Applied Mathematics and Statistics;
Department of Computer Science, Stony Brook University\\[2pt]
\textbf{Ver\"offentlicht:} ICLR 2025 (OpenReview: J9FgrqOOni)\\[2pt]
\textbf{Code:} \url{https://github.com/wenhangao21/ICLR25-CROP}

% =========================================================
\subsection*{Motivation und Problemstellung}
% =========================================================

Eine der zentralen Eigenschaften, mit denen der Fourier Neural Operator (FNO)
beworben wird, ist die sogenannte \emph{Diskretisierungsinvarianz}: Da die
gelernten Kerneloperatoren im Funktionenraum definiert sind und die Netzwerkparameter von der Gitteraufl\"osung entkoppelt sind, soll das Modell prinzipiell
auf Aufl\"osungen angewendet werden k\"onnen, die von der Trainingsaufl\"osung
abweichen. Diese Eigenschaft hat in der Praxis dazu gef\"uhrt, dass FNO-Modelle
routinem\"a\ss ig auf einer niedrigeren Aufl\"osung trainiert und auf einer
h\"oheren getestet werden (\emph{zero-shot super-resolution}).

Die Autoren identifizieren dies als ein verbreitetes Missverst\"andnis in der
Operator-Learning-Community: W\"ahrend die Parameter des FNO tats\"achlich
aufl\"osungsunabh\"angig sind und das Modell Eingaben verschiedener Aufl\"osungen
verarbeiten \emph{kann}, garantiert dies keineswegs, dass die Ausgabequalit\"at
konsistent bleibt. Ziel der Arbeit ist es, dieses Missverst\"andnis formell
zu adressieren, die Fehlerquelle zu isolieren und eine L\"osungspipeline
vorzuschlagen.

% =========================================================
\subsection*{Kernbeitrag: Discretization Mismatch Error (DME)}
% =========================================================

Die Autoren f\"uhren den Begriff des \textbf{Discretization Mismatch Error
(DME)} ein, definiert als die Diskrepanz zwischen den FNO-Ausgaben, wenn
dieselbe kontinuierliche Eingabefunktion in zwei verschiedenen Aufl\"osungen
diskretisiert wird:
\[
  E_{MN} := \bigl\| \mathcal{G}_\theta\bigl(a|_{\Omega_M}\bigr)
             - \mathcal{G}_\theta\bigl(a|_{\Omega_N}\bigr) \bigr\|_{\mathcal{U}}.
\]
Sie unterscheiden bewusst zwischen der klassischen Definition von
Diskretisierungsinvarianz nach Kovachki et al.\ (2021) -- Konvergenz des
diskreten Operators gegen den kontinuierlichen Grenzoperator -- und der
praktisch relevanten Frage, ob die Modellausgaben bei verschiedenen
Diskretisierungen \emph{konsistent} sind.

\textbf{Proposition 4.4} zeigt theoretisch, dass der DME mit wachsender
Anzahl von Fourier-Schichten akkumuliert und mit zunehmendem Abstand
zwischen Trainings- und Inferenzaufl\"osung w\"achst. Der Fehler ist durch
einen Ausdruck beschr\"ankt, der von den Lipschitz-Konstanten der
Aktivierungsfunktionen, den Operatornormen der linearen Schichten und
dem relativen Aufl\"osungsunterschied $|M - N|/(MN)$ abh\"angt.

Zus\"atzlich zu DMEs identifizieren die Autoren \emph{Aliasing-Fehler} als
verwandte, aber separate Fehlerquelle -- bereits von Raoni\v{c} et al.\
(CNO, NeurIPS 2023) diskutiert -- und zeigen, dass alias-freie
Aktivierungsfunktionen die DME-Problematik nicht beheben.

% =========================================================
\subsection*{Vorgeschlagene L\"osung: CROP-Pipeline}
% =========================================================

Als Reaktion auf die DME-Analyse schlagen die Autoren die
\textbf{Cross-Resolution Operator-learning Pipeline (CROP)} vor. CROP
besteht aus zwei Bausteinen:

\begin{itemize}
  \item \textbf{CROP Lifting Layer:} Bildet die Eingabefunktion durch eine
    globale bandlimitierte Faltung auf eine latente Repr\"asentation mit
    fester Diskretisierung ab. Da die Bandlimitierung die Diskretisierung
    des latenten Raums fixiert, entstehen keine DMEs.

  \item \textbf{CROP Projection Layer:} Bildet die latente Funktion durch
    eine bandlimitierte Faltung kombiniert mit einem punktweisen
    vollverbundenen Netz auf den Ausgaberaum ab. Das vollverbundene Netz
    kann dabei Frequenzanteile \"uber die Bandgrenze hinaus erzeugen,
    sodass auch feine Details rekonstruiert werden k\"onnen.
\end{itemize}

Der Zwischenoperator zwischen Lifting und Projection kann frei gew\"ahlt
werden -- sowohl FNO als auch U-Net sind m\"oglich -- was CROP zu einer
generischen Pipeline macht. Die CROP-Architektur erf\"ullt
Continuous-Discrete Equivalence (CDE) im Sinne von Raoni\v{c} et al.\
(2023) und ist somit frei von Aliasing- und Diskretisierungsfehlern.

% =========================================================
\subsection*{Experimentelle Ergebnisse}
% =========================================================

Die Autoren evaluieren CROP auf der 2D-Navier-Stokes-Gleichung sowie
Darcy-Flow, Helmholtz und Poisson, mit besonderem Fokus auf
Cross-Resolution-Aufgaben.

\begin{itemize}
  \item \textbf{DME-Akkumulation (Navier-Stokes, Re = 1000):} Ein auf
    $64 \times 64$ trainierter FNO erzielt bei Inferenz auf
    $128 \times 128$ einen dreimal h\"oheren Fehler als bei der
    Trainingsaufl\"osung (0{,}66\,\% vs.\ 0{,}22\,\% relativer
    $\ell_2$-Fehler, bereits im ersten Zeitschritt). Der DME w\"achst
    monoton mit der Anzahl der Fourier-Schichten.

  \item \textbf{Cross-Resolution (Navier-Stokes):} FNO zeigt starke
    Abh\"angigkeit von der Trainingsaufl\"osung (Fehler 0{,}58\,\%
    bei $64\times64$, 4{,}65\,\% bei $256\times256$). CROP erreicht
    konsistent 0{,}54\,\% auf allen getesteten Aufl\"osungen
    ($32\times32$ bis $256\times256$).

  \item \textbf{Lernf\"ahigkeit (High-Re Navier-Stokes, Darcy, Helmholtz,
    Poisson):} CROP \"ubertrifft alle Baselines (FNO, CNO, U-Net, ResNet,
    DeepONet) auch bei der In-distribution-Genauigkeit, da die
    bandlimitierten Lifting- und Projection-Schichten globale und lokale
    Merkmale effizient kombinieren.
\end{itemize}

% =========================================================
\subsection*{Bezug zur Masterarbeit}
% =========================================================

\subsubsection*{1.\ Direkte Rechtfertigung der Einschr\"ankung im Related Work}

Die Arbeit liefert die formelle theoretische Grundlage f\"ur den Hinweis
in Kapitel~2.3 der Masterarbeit, dass die Diskretisierungsinvarianz des FNO
in der Praxis nur n\"aherungsweise gilt. Der DME-Begriff erkl\"art pr\"azise,
warum: Die Netzwerkparameter sind zwar aufl\"osungsunabh\"angig, aber die
diskrete Implementierung via FFT akkumuliert Fehler, sobald Trainings- und
Inferenzaufl\"osung voneinander abweichen.

\subsubsection*{2.\ Kein direkter Einfluss auf die Masterarbeit-Experimente}

Da alle Experimente der Masterarbeit auf einer festen $256 \times 256$-Aufl\"osung
durchgef\"uhrt werden -- sowohl Training als auch Inferenz -- tritt der DME
in der vorliegenden Arbeit nicht auf. Die Relevanz der Gao et al.-Arbeit liegt
daher prim\"ar in der \emph{konzeptionellen Einordnung}: Die Behauptung aus dem
FNO-Paper, das Modell sei diskretisierungsinvariant, ist mit Vorsicht zu
genie\ss en und sollte nicht unkommentiert \"ubernommen werden.

\subsubsection*{3.\ Spektrale Bias-Problematik}

Ein Nebenbefund der Gao et al.-Arbeit ist, dass FNO auch bei gleicher
Trainings- und Inferenzaufl\"osung zur Gl\"attung feiner Details neigt, da
hohe Fourier-Moden abgeschnitten werden. Dies ist konsistent mit den in der
Masterarbeit beobachteten Residuen an Kristallgrenzen, wo der FNO steile
Gradienten in $\psi$ unterschätzt.

\subsubsection*{4.\ CROP als m\"ogliche Weiterentwicklung}

Die CROP-Architektur k\"onnte als Ausgangspunkt f\"ur eine Weiterentwicklung
des FNO-Surrogatmodells dienen, insbesondere wenn k\"unftige Arbeiten
verschiedene Gitteraufl\"osungen untersuchen oder den Fehler an
Kristallgrenzen durch bessere Hochfrequenzrepresentation reduzieren wollen.

% =========================================================
\subsection*{Kritische Einsch\"atzung}
% =========================================================

\begin{itemize}
  \item Die Arbeit konzentriert sich auf Standard-PDE-Benchmarks (Navier-Stokes,
    Darcy, Poisson, Helmholtz) ohne starre Einschl\"usse oder variable
    Geometrien. \"Ubertragbarkeit auf geometrisch komplexere Probleme wie
    Mehrkristall-Stokes-Str\"omung ist nicht direkt demonstriert.

  \item Die CROP-Pipeline erfordert zus\"atzliche Hyperparameter (Bandlimit,
    Wahl des Zwischenoperators), was den Abstimmungsaufwand gegen\"uber einem
    einfachen FNO erh\"oht.

  \item Der DME ist konzeptionell klar, aber die vorgestellten Schranken
    (Proposition 4.4) sind obere Schranken und k\"onnen in der Praxis sehr
    konservativ sein -- die tats\"achliche Fehlerakkumulation h\"angt stark
    vom Problem ab.
\end{itemize}

\end{document}
